\subsection{The precise ideal class paired with the norm target}

\begin{statement}{A.155}{Proved}
{A class-group criterion for the actual fixed-support norm family}
Let
\[
 a=3\cdot11^{M_1}59^{M_2},\qquad 24\mid M_1,M_2>0,
 \qquad K_n=\Q(\alpha_n),\quad \alpha_n^{3^n}=a,
 \quad\alpha_{n+1}^3=\alpha_n.
\tag{A.155.1}
\]
Write \(S=\{2,3,11,59,\infty\}\) and let \(S_n\) be the places
of \(K_n\) above \(S\).
For \(v=2,11,59\), choose the compatible residue-degree-one branch
\(P_{v,0,n}\) described in the proof.  For any prime \(P\) above
\(v\), put \(f_P=[\kappa(P):\F_v]\).
Define the finite three-primary group
\[
 C_n^{\rm bal}=
 \operatorname{Cl}(K_n)\{3\}\big/
 \left\langle
   [P]_{(3)}-f_P[P_{v,0,n}]_{(3)}:
       v\in\{2,11,59\},\ P\mid v
 \right\rangle,
 \qquad e_n=[P_{2,0,n}]\in C_n^{\rm bal}.
\tag{A.155.2}
\]
Here \([P]_{(3)}\) is the projection of the ordinary ideal class
to its three-primary summand; primes are not inverted in that class
group before the indicated quotient is taken.

For \(R_r=\Z/3^r\), let
\[
 \begin{gathered}
 P_{n,r}=R_r[\operatorname{Hom}_{\Q}(K_n,\overline\Q)],\\
 I_{n,r}=\ker(P_{n,r}\xrightarrow{\sum}R_r),\qquad
 J_{n,r}=P_{n,r}/R_r\mathbf1,\qquad
 M_{n,r}=I_{n,r}\otimes\mu_{3^r}.
 \end{gathered}
\tag{A.155.3}
\]
The actual Cartier dual of \(M_{n,r}\) is \(J_{n,r}\).
There is a canonical isomorphism, specified by the unramified
normalization in the proof,
\[
 \Sha^1_S(J_{n,r})
    :=\ker\left(H^1(G_{\Q,S},J_{n,r})
               \longrightarrow\prod_{v\in S}H^1(\Q_v,J_{n,r})\right)
       \simeq\operatorname{Hom}(C_n^{\rm bal},R_r).
\tag{A.155.4}
\]
Let \(o_{n,r}\) be the image of \(\kappa_{3^r}(-2)\) under
the connecting map of
\[
 0\longrightarrow M_{n,r}\longrightarrow
       P_{n,r}\otimes\mu_{3^r}
         \xrightarrow{\sum}\mu_{3^r}\longrightarrow0.
\tag{A.155.5}
\]
It is locally zero at every place in \(S\).
Using arithmetic Frobenius and the cup-product convention made
explicit below, its Poitou--Tate pairing is exactly
\[
                    \langle c,o_{n,r}\rangle
                         =\frac{h(e_n)}{3^r}\quad\text{in }\Q/\Z,
 \qquad c\longleftrightarrow h
       \text{ under (A.155.4)}.
\tag{A.155.6}
\]

Consequently the following two assertions are equivalent:
\[
 \begin{aligned}
 &-2\in N_{K_n/\Q}
           \widehat{\mathcal O_{K_n,S_n}^{\times}}_{(3)},\\
 &e_n=0\quad\text{in }C_n^{\rm bal}.
 \end{aligned}
\tag{A.155.7}
\]
Ideal norm gives actual maps
\(C_{n+1}^{\rm bal}\to C_n^{\rm bal}\) taking
\(e_{n+1}\) to \(e_n\).
There exists one family
\[
 b_n\in\widehat{\mathcal O_{K_n,S_n}^{\times}}_{(3)},\qquad
 N_{K_{n+1}/K_n}b_{n+1}=b_n,\qquad N_{K_n/\Q}b_n=-2
\tag{A.155.8}
\]
if and only if \(e_n=0\) for every \(n\).
Moreover \(3^n e_n=0\).
The tests in (A.155.6) retain all coefficient precisions: this order
bound on the element does not bound the precision required to detect
it inside the ambient class group.
The theorem identifies the exact class and pairing to compute.
It does not assert that these classes vanish in the two-prime tower.
\end{statement}

\begin{proof}
\emph{The local branches and the local norm target.}
At two, the cube map is an automorphism of \(\Q_2^\times\)
on the unit group.  The odd unit \(a\) therefore has a compatible
root of every three-power order in \(\Q_2\), giving the distinguished
degree-one branches \(P_{2,0,n}\).
At \(q=11,59\), write \(M_q=3^s m\) with \(3\nmid m\).
Every \(q\)-adic unit has compatible three-power roots, since
\(3\nmid q-1\).
For \(n\leq s\), take the resulting degree-one factor of the
root algebra.  For \(n>s\), take its factor generated by a compatible
root of \(q\) of degree \(3^{n-s}\).
The Bezout root extraction in A.151 identifies this factor literally
with the chosen factor of \(K_n\otimes\Q_q\).
Its polynomial can be chosen Eisenstein, so its residue degree is
one.  These choices are compatible as \(n\) increases.

The radicand has three-adic valuation one, so the unique three-prime
of \(K_n\) is totally ramified and has residue degree one.
The local and global splitting fields, even after adjoining extra
three-power roots of unity, have their full holomorph groups, by the
Eisenstein and non-Galois cubic argument in A.149.
In particular the pure local field at three has no nontrivial
abelian subextension of \(\Q_3\); its compositum with every finite
three-power cyclotomic extension is totally ramified.

The target \(-2\) is a norm from each local root algebra.
At two put it on the degree-one factor and put one on the others.
At \(q=11,59\) it is a unit: on the selected factor of degree
\(3^j\), choose its scalar \(3^j\)-th root in \(\Q_q^\times\).
The norm of that scalar is exactly \(-2\).
At three use the proved local norm surjectivity A.140.
At the real place use the real factor of the odd-degree polynomial.
These facts prove the local vanishing claimed after (A.155.5).

\emph{Lifting a locally trivial dual class to a class character.}
Use the literal dual coefficient sequence
\[
              0\longrightarrow R_r\mathbf1\longrightarrow
                    P_{n,r}\longrightarrow J_{n,r}\longrightarrow0.
\tag{A.155.9}
\]
The group \(\Sha^1_S(\mu_{3^r})\) is zero.
Indeed a Kummer class in \(H^1(G_{\Q,S},\mu_{3^r})\) is represented
by a rational \(S\)-unit.  Local triviality at the finite places of
\(S\) makes all its valuations divisible by \(3^r\).
Dividing by a rational \(3^r\)-th power leaves a sign, which itself
is a \(3^r\)-th power.
Finite duality \cite[I, Theorem 4.10(a)]{BIR_MilneADT} consequently
gives \(\Sha^2_S(R_r)=0\).
The connecting class of any \(c\in\Sha^1_S(J_{n,r})\) in
(A.155.9) is in this zero group.
Continuous Shapiro therefore gives a lift
\[
       \widetilde h\in H^1(G_{\Q,S},P_{n,r})
                  =H^1(G_{K_n,S_n},R_r).
\tag{A.155.10}
\]
For every \(v\in S\), local exactness says that the restrictions
of \(\widetilde h\) on all factors above \(v\) are the restrictions
of one character \(\chi_v\in H^1(\Q_v,R_r)\).

At \(v=2,11,59\), every such character is unramified: the pro-three
part of the local unit group is zero, so local reciprocity leaves only
the valuation character.
At three, the restriction of global characters
\(H^1(G_{\Q,S},R_r)\) onto the ramified characters of
\(\Q_3^\times\) is an isomorphism.
To see this exact assertion, local reciprocity kills the inertia at
\(2,11,59\) for every cyclic three-power character.
The ray class sequence for \(\Q\) thus leaves precisely the
three-primary cyclotomic characters ramified at three.
Their restriction to \(1+3\Z_3\) is every continuous character
of that group with values in \(R_r\); the sign has order two.
This uses the reciprocity and existence maps of
\cite[V, Theorems 1.7 and 3.5--3.6]{MilneCFT}.

Subtract the unique global character with the same three-adic
inertia restriction as \(\chi_3\).
The resulting lift \(h\) in (A.155.10) is unramified at every finite
place of \(K_n\), and hence is an ordinary ideal-class character
\(\operatorname{Cl}(K_n)\{3\}\to R_r\).
At each \(v=2,11,59\), restriction of an unramified base character
sends arithmetic Frobenius at \(P\) to its \(f_P\)-th power.
Thus the exact conditions on this character are
\[
                    h([P])=f_P h([P_{v,0,n}]).
\tag{A.155.11}
\]
At three there is only one prime, with residue degree one, so any
unramified character there is the restriction of an unramified base
character.  There is no additional relation.
This gives the map in (A.155.4) into the stated character group.

Conversely, any ideal-class character satisfying (A.155.11) is
locally the restriction of a base character at every place of \(S\).
Its Shapiro class projects to a class in \(\Sha^1_S(J_{n,r})\).
Finally, an unramified class character in the kernel of this projection
is the restriction of a global three-primary cyclotomic character.
If that character is nonzero, its restriction to the pure local
three-field is still nonzero and ramified: the full local holomorph
has cyclotomic abelianization, and the compositum is totally ramified,
as proved above.  This contradicts unramifiedness.
The kernel is zero, proving (A.155.4) and uniqueness of its normalization.

\emph{The explicit cochain pairing and its sign.}
Write \(b=\kappa_{3^r}(-2)\), and choose a continuous one-cochain
\(B\) in \(P_{n,r}\otimes\mu_{3^r}\) lifting \(b\).
Then \(dB\) represents \(o_{n,r}\).
Represent the normalized lift in (A.155.10) by a cocycle \(h\)
in \(P_{n,r}\), and let \(c\) be its image in \(J_{n,r}\).
At each place choose a zero-cochain \(w_v\) in \(P_{n,r}\) such that
\[
                         h_v=\chi_v\mathbf1+dw_v.
\tag{A.155.12}
\]
All pairings here are the coordinate dot products.
Because \(dh=0\), the global cochain \(-h\cup B\) has differential
\(h\cup dB\), which is the pairing of \(c\) and \(o_{n,r}\).
The local null-homotopy is \(w_v\cup dB_v\).
The local-minus-global convention for the Poitou--Tate pairing,
as in the explicit description following
\cite[I, Remark 4.19]{BIR_MilneADT}, therefore gives the local class
\[
 \begin{aligned}
       w_v\cup dB_v+h_v\cup B_v
          &=\chi_v\cup b_v+d(w_v\cup B_v),\\
       \langle c,o_{n,r}\rangle
          &=\sum_{v\in S}\operatorname{inv}_v(\chi_v\cup b_v).
 \end{aligned}
\tag{A.155.13}
\]
This calculation retains the actual coefficient morphisms of
(A.155.5) and (A.155.9).
For the order used here, local reciprocity identifies
\(\operatorname{inv}(\chi\cup\kappa(b))\) with
\(\chi(\operatorname{Art}(b))/3^r\);
see \cite[III, Proposition 3.6 and Section 4, Steps 2--3]{MilneCFT}.
Every normalized \(\chi_v\) is unramified.
The rational number \(-2\) has valuation one at two and zero at
the other finite places in \(S\).
Consequently (A.155.13) is exactly (A.155.6).

\emph{From all finite coefficient tests to one norm fibre.}
Finite Poitou--Tate duality is perfect on the two localization kernels.
Thus (A.155.6) vanishes for every \(r\) and every \(h\) if and only
if \(e_n=0\) in the finite group (A.155.2).
Let \(I_n\) be the integral augmentation kernel on the \(3^n\)
root labels, and put \(M_n=I_n\otimes\Z_3(1)\).
The finite groups \(H^1(G_{\Q,S},M_n/3^r)\) are finite by
\cite[I, Corollary 4.15]{BIR_MilneADT}.
Coefficient transitions lift continuous cochains, and the resulting
\(\varprojlim^1H^1\) vanishes by the finite Mittag--Leffler property.
Hence
\[
 H^2_{\rm cts}(G_{\Q,S},M_n)
            =\varprojlim_r H^2(G_{\Q,S},M_n/3^r).
\tag{A.155.14}
\]
The integral connecting class of \(\kappa_3(-2)\) is zero exactly
when all \(o_{n,r}\) are zero.
The integral permutation module is continuously split over its
augmentation as an underlying \(\Z_3\)-module.
Shapiro and the integral Kummer sequence identify its degree-one
cohomology with
\(\widehat{\mathcal O_{K_n,S_n}^{\times}}_{(3)}\): the inverse Tate
module of the finite \(S_n\)-class group is zero.
Under this identification augmentation is the actual field norm.
The integral long exact sequence proves (A.155.7).

\emph{The literal tower transitions and compact norm fibres.}
For \(P'\) over \(P\), write \(f_{P'/P}\) for its relative
residue degree.  The chosen degree-one primes have relative residue
degree one, and ideal norm sends them to the preceding chosen primes.
Moreover
\[
 \begin{aligned}
 N\bigl([P']-f_{P'}[P_{v,0,n+1}]\bigr)
       &=f_{P'/P}\bigl([P]-f_P[P_{v,0,n}]\bigr),\\
 N e_{n+1}&=e_n.
 \end{aligned}
\tag{A.155.15}
\]
This proves the asserted maps on the actual quotients.
On root labels the primal coefficient map is
\(e_j\mapsto e_{j\bmod3^n}\); its Cartier dual is
\(e_i\mapsto\sum_{j\equiv i\bmod3^n}e_j\).
Shapiro identifies the latter map on class characters with
\(h\mapsto h\circ N_{K_{n+1}/K_n}\), so these are the same
transitions as (A.155.15).
At two the extension is unramified and
\(\sum_{P\mid2}f_P=3^n\).
The principal ideal \((2)=\prod_{P\mid2}P\), followed by
(A.155.2), therefore gives \(3^n e_n=0\).
This is an order bound on \(e_n\), not a coefficient-detection bound.
For example, the element three in \(\Z/9\) has order three
but is killed by every homomorphism to \(\Z/3\).
The equivalence above uses all \(r\), or a precision at least
the exponent of the entire finite group \(C_n^{\rm bal}\).
Each completed \(S_n\)-unit group is compact, by the finite
generation of \(S_n\)-units.
Its fibre over \(-2\) under norm to \(\Q\) is closed and therefore
compact.  If every \(e_n=0\), (A.155.7) makes all these fibres
nonempty.  A point in any largest finite height gives all earlier
points by the actual field norms.  The compact finite-intersection
argument gives (A.155.8), with no assumption that individual
transition maps on the fibres are surjective.
The converse follows by projection to each level.
\end{proof}

The finite checks in
\path{code/balanced_class_norm_criterion_certificate.py}
retain the cochain differential identity, the residue-degree norm
relations, and the exact finite-character separation used above.
They do not calculate \(e_n\) in the actual class groups.
